\section{書式と例}
\subsection{単著・共著（\bibtype{book}）}
著者が1人の書籍は，\bibtype{book} に書きます。和文の著者には，並べ替え用の \bibfield{sortname}（ローマ字）を添えます。

記入のひな形は次のとおりです：
\begin{bibtemplate}
@book{<キー>,
  author    = {<姓>,<名>},         ※ 和文は姓と名を「,」で区切る
  sortname  = {<Family>, <Given>}, ※ 和文のみ。ローマ字
  date      = {<刊行年>},
  title     = {<書名>},
  publisher = {<出版社>},
  langid    = {japanese},
}
\end{bibtemplate}
たとえば，次のように書きます：
\bibexample{aoki2015}
共著は，著者を \texttt{and} でつなぎます。文献リストでは中黒でつながります：
\bibexample{aokikita2018}
欧文の共著は，2人目以降を名・姓の順に出し，最後を and でつなぎます：
\bibexample{clark2017}

\subsection{翻訳書（\bibfield{related}）}
訳書は原典のエントリを作り，\bibfield{related} でつなぎます。文献リストには，原典の項目に訳書が添えられて1項目で出ます：
\bibexample{hale1962,hale1985}
原典を引用していなくても，原典のエントリを bib に書いておけば，同じ形で出ます（\bibfield{related} が指す原典が代表になります）：
\bibexample{hale1985}

\subsection{LaTeX の例}
LaTeX の例は，薄い青の地で示します：
\begin{ltxexample}
\documentclass{jlreq}
\usepackage[backend=biber,style=jsr]{biblatex}
\addbibresource{refs.bib}
\end{ltxexample}

\section{コマンドとフィールド}
\begin{ltxsyntax}
\cmditem{parencite}[prenote][postnote]{key}
\cmditem{textcite}[prenote][postnote]{key}

全角の丸括弧で囲む文献注です。本文での引用の基本形です。ページは \prm{postnote} に書きます。

\cmditem{citeauthor}[prenote][postnote]{key}

著者名だけを出します。
\end{ltxsyntax}
\begin{fieldlist}
\fielditem{origpublisher}{literal}

紙媒体の出版社。電子書籍で使います。

二つ目の段落です。段落頭は下げず，段落の間を少し空けます。
\listitem{location}{literal}

出版地。欧文だけ出します。
\end{fieldlist}
\begin{typelist}
\typeitem{book}

単著・共著の書籍です。

\reqitem{author, title, date}
\optitem{subtitle, edition, publisher, location}
\end{typelist}
\begin{optionlist}
\boolitem[true]{firstfull}

文献注で，初出の人物をフルネームで出します。
\optitem[false]{doi}{\opt{true}, \opt{false}}

取得日がない文献にも DOI を出します。
\end{optionlist}
\begin{keymarglist}
\item[translationof] 訳書。原典の項目に添えます。
\item[reprintof] 再録・再刊。
\end{keymarglist}
